\section{Introduction}
\label{sec:intro}
Long-context model execution makes the memory hierarchy difficult to ignore. An implementation can perform mathematically familiar operations yet exhibit very different runtime behavior depending on which tensors are stored, when they are reused, and how much work is exposed to the device. Describing a workload only by its floating-point operation count therefore leaves an important part of the systems problem unspecified. The missing description is a schedule: which intermediate values are produced together, which remain nearby, and which travel through a more expensive level of memory.

The Transformer architecture established attention as a central sequence-modeling operation~\cite{vaswani2017attention}. Subsequent IO-aware work reorganized attention computation around tiling and reuse, rather than treating the memory traffic of a straightforward schedule as unavoidable~\cite{dao2022flashattention}. These developments motivate a practical question for readers of performance papers: what evidence distinguishes a reduction in traffic, a reduction in peak allocation, and an improvement in end-to-end time? The three quantities are related but cannot be substituted for one another.

This manuscript addresses that question through an explicitly synthetic study. Its purpose is not to introduce a new attention kernel. It provides a small model whose assumptions can be inspected, whose figures can be regenerated, and whose limitations are visible next to its conclusions. The model assigns arithmetic and byte counts to two abstract schedules and combines them with hypothetical device parameters. The resulting curves are numerical consequences of those assignments. They are not estimates calibrated to a particular accelerator, and they are not reproductions of published benchmark results.

The distinction matters because a convincing-looking graph can conceal several different claims. A smooth increase in speedup might result from improved reuse, but it might also follow automatically from a chosen denominator. A low workspace curve might exclude an allocation that a real runtime cannot avoid. A threshold at which one schedule becomes faster might move when launch latency, batching, or synchronization is included. Keeping the equations and the plotting data in the same project makes these possibilities easier to examine.

\subsection{Questions and intended use}
We organize the study around four questions. First, how does an explicit quadratic intermediate affect modeled memory traffic and workspace as sequence length grows? Second, when does reducing that traffic stop helping because arithmetic becomes the active constraint? Third, which assumptions move the crossover between the two abstract schedules? Fourth, what experimental protocol would be necessary to turn the illustration into a defensible implementation study?

A reader can use the artifact at three levels. At the first level, the figures offer a visual explanation of data movement. At the second, the equations expose a compact performance model that can be challenged or replaced. At the third, the project is a working authoring example: citations are resolved by BibTeX, figures are genuine vector graphics, and the preview is generated by a real LaTeX engine. None of these software properties establishes the truth of a hardware claim; they make the artifact easier to audit.

The remainder of the paper separates these levels deliberately. Section~\ref{sec:background} introduces the computational context. Sections~\ref{sec:model} and~\ref{sec:design} specify the model and its artifact organization. Sections~\ref{sec:evaluation} and~\ref{sec:sensitivity} interpret the synthetic results. Section~\ref{sec:discussion} explains how a future measured study would need to differ. The appendix records additional checks and reporting decisions so that the example can be extended without losing the distinction between evidence and illustration.
